\documentclass[11pt,a4paper]{article}
\usepackage[margin=2.1cm]{geometry}
\usepackage{lmodern}
\usepackage[T1]{fontenc}
\usepackage{amsmath,amssymb,bm,mathtools}
\usepackage{booktabs}
\usepackage{enumitem}
\usepackage{microtype}
\setlength{\parindent}{0pt}\setlength{\parskip}{6pt plus 2pt minus 1pt}
\allowdisplaybreaks

\newcommand{\vect}[1]{\bm{#1}}
\newcommand{\mat}[1]{\bm{#1}}
\newcommand{\T}{^{\mathsf{T}}}
\newcommand{\Real}{\mathbb{R}}
\newcommand{\diag}{\operatorname{diag}}
\DeclareMathOperator{\blkdiag}{blkdiag}
\newcommand{\Rot}{\mat{R}}
\newcommand{\hx}{\hat{\vect{x}}}

\title{Docking EKF: equation set}
\author{}
\date{8 October 2026}

\begin{document}
\maketitle
\vspace{-2em}

\section{State, inputs and timing}
\begin{equation}
\vect x = \begin{bmatrix} p_x \\ p_y \\ \theta \\ v_x \\ v_y \end{bmatrix}\in\Real^5,
\qquad
\vect u = \begin{bmatrix} a_x \\ a_y \\ r \end{bmatrix}\in\Real^3,
\qquad
T = \tfrac{1}{200}\ \text{s}.
\end{equation}
$\vect p$ and $\vect v$ are the position and velocity of the chassis centre expressed in the target frame $T$; $\theta$ is the heading of the body $x$-axis relative to the target $x$-axis. $a_x,a_y$ are the gravity-compensated body-frame accelerations and $r$ the yaw rate from the IMU. The filter steps at the IMU period $T$.

\section{Continuous model}
\begin{equation}
\dot p_x = v_x,\qquad
\dot p_y = v_y,\qquad
\dot\theta = r,\qquad
\dot v_x = \cos\theta\,a_x - \sin\theta\,a_y,\qquad
\dot v_y = \sin\theta\,a_x + \cos\theta\,a_y.
\end{equation}

\section{Discrete map (Euler)}
\begin{equation}
\vect x_{k+1} = \vect f_d(\vect x_k,\vect u_k) = \vect x_k + T
\begin{bmatrix}
v_x \\ v_y \\ r \\ \cos\theta\,a_x - \sin\theta\,a_y \\ \sin\theta\,a_x + \cos\theta\,a_y
\end{bmatrix},
\qquad \theta \text{ wrapped to } [-\pi,\pi).
\end{equation}

\section{Jacobians of the discrete map}
Evaluated at the estimate before propagation and at the current IMU sample:
\begin{equation}
\mat F_k = \frac{\partial\vect f_d}{\partial\vect x} =
\begin{bmatrix}
1 & 0 & 0 & T & 0 \\
0 & 1 & 0 & 0 & T \\
0 & 0 & 1 & 0 & 0 \\
0 & 0 & -T(\sin\theta\,a_x + \cos\theta\,a_y) & 1 & 0 \\
0 & 0 & \phantom{-}T(\cos\theta\,a_x - \sin\theta\,a_y) & 0 & 1
\end{bmatrix},
\qquad
\mat G_k = \frac{\partial\vect f_d}{\partial\vect u} =
\begin{bmatrix}
0 & 0 & 0 \\
0 & 0 & 0 \\
0 & 0 & T \\
T\cos\theta & -T\sin\theta & 0 \\
T\sin\theta & \phantom{-}T\cos\theta & 0
\end{bmatrix}.
\end{equation}

\section{Process noise}
\begin{gather}
\mat\Sigma_u = q_{\mathrm{scale}}\,\diag(\sigma_a^2,\sigma_a^2,\sigma_r^2),
\qquad
\mat Q_k = \mat G_k\mat\Sigma_u\mat G_k\T + \mat Q_{\mathrm{floor}},\\
\mat Q_{\mathrm{floor}} = \diag\!\big(\epsilon,\ \epsilon,\ q_\theta,\ q_v,\ q_v\big).
\end{gather}
$\sigma_a,\sigma_r$ are the per-sample IMU white-noise standard deviations, inflated by $q_{\mathrm{scale}}$. $q_v$ and $q_\theta$ are the per-step process-noise variances on velocity and heading: tuning parameters, chosen so that the filter is consistent (time-averaged NEES near 5 with the marker visible). They express how far the velocity and heading predictions may drift from the IMU integration between fixes, whether through residual bias, slip or target motion.

\section{Predictor (every IMU sample)}
\begin{equation}
\hx_k^- = \vect f_d(\hx_{k-1},\vect u_k),
\qquad
\mat P_k^- = \mat F_k\mat P_{k-1}\mat F_k\T + \mat Q_k.
\end{equation}

\section{Fiducial measurement (15 Hz, captured at $t_c = t_k - \tau_d$)}
Pre-transform of the raw solver output, the planar translation $\vect t$ and yaw $\psi$ of the marker in the camera frame, with $\vect m$ the marker position in $T$ and $\vect l_c$ the camera position in the body frame:
\begin{equation}
\theta_z = \psi_0 - \psi,
\qquad
\vect p_z = \vect m - \Rot(\theta_z)\,(\vect t + \vect l_c).
\end{equation}
The pseudo-measurement is then linear in the state:
\begin{gather}
\vect z = \begin{bmatrix}\vect p_z \\ \theta_z\end{bmatrix} = \mat H_f\vect x + \vect v,
\qquad
\mat H_f = \begin{bmatrix} 1 & 0 & 0 & 0 & 0 \\ 0 & 1 & 0 & 0 & 0 \\ 0 & 0 & 1 & 0 & 0 \end{bmatrix},\\
\mat R_f = \diag(\sigma_p^2,\sigma_p^2,\sigma_\theta^2),
\qquad
\sigma_p = \sigma_{p,0}\sqrt{1 + (d/d_0)^2}\quad\text{with } d \text{ the range to the marker}.
\end{gather}

\section{Corrector (when the fix becomes available), with delay compensation}
\begin{align}
\vect x_{\mathrm{ref}} &= \text{buffer entry at the capture time } t_c \quad(\text{or } \hx_k^- \text{ if compensation is off}), \\
\vect\nu &= \vect z - \mat H_f\vect x_{\mathrm{ref}}, \qquad \nu_3 \text{ wrapped}, \\
\mat S &= \mat H_f\mat P_k^-\mat H_f\T + \mat R_f, \\
\mat K &= \mat P_k^-\mat H_f\T\mat S^{-1}, \\
\hx_k &= \hx_k^- + \mat K\vect\nu, \qquad \theta \text{ wrapped}, \\
\mat P_k &= (\mat I - \mat K\mat H_f)\,\mat P_k^-\,(\mat I - \mat K\mat H_f)\T + \mat K\mat R_f\mat K\T \qquad\text{(Joseph form)}.
\end{align}

\section{Fiducial outlier gate}
Each fix is accepted only if $\vect\nu\T\mat S^{-1}\vect\nu \le 100$ (about ten standard deviations); otherwise it is discarded and the state and covariance are left unchanged. The threshold is deliberately loose: it rejects a wrong-marker or degenerate pose solve, metres off, but never a plausible fix, so an unmodelled change such as a target starting to move is still absorbed rather than rejected.

RTK position and INS heading are not used by the docking filter. In the navigation filter, which works in the world frame, they are the rows
$\vect z_W = [\vect p_{\mathrm{gps}};\ \psi_W]$, $\mat H = [\mat I_3\ \ \mat 0_{3\times 2}]$, $\mat R = \diag(\sigma_{\mathrm{rtk}}^2,\sigma_{\mathrm{rtk}}^2,\sigma_\psi^2)$, heading innovation wrapped.

\section{Initialisation}
First fiducial fix, velocity unknown (deployment form, both variants):
\begin{equation}
\hx_0 = \begin{bmatrix}\vect z_f \\ \vect 0\end{bmatrix},
\qquad
\mat P_0 = \blkdiag\!\big(\mat R_f,\ \sigma_{v,0}^2\mat I_2\big).
\end{equation}
Handover from the navigation estimate (full-scale alternative, needs the target's reported pose):
\begin{gather}
\hx_0 = \begin{bmatrix}\Rot(\theta_t)\T(\hat{\vect p}_W - \vect p_{t,W}) \\ \hat\theta_W - \theta_t \\ \Rot(\theta_t)\T\hat{\vect v}_W\end{bmatrix},
\qquad
\mat T_m = \blkdiag\!\big(\Rot(\theta_t)\T,\ 1,\ \Rot(\theta_t)\T\big),\\
\mat P_0 = \mat T_m\mat P_W\mat T_m\T + \blkdiag(\mat P_{\mathrm{tgt}},\mat 0).
\end{gather}

\section{Schedule, per IMU sample $k$}
\begin{enumerate}[nosep]
  \item Predict (Section~6), unless $k = 1$.
  \item For every fiducial fix with $t_{\mathrm{avail}} \le t_k$ not yet used: if valid, gate (Section~9) and correct (Section~8).
  \item Store $\hx_k$ in the ring buffer at index $k$, \emph{after} the corrections, so that a later delayed fix sees the best estimate of time $t_k$.
  \item Log $\hx_k$ and $\diag(\mat P_k)$. With truth available, $\mathrm{NEES}_k = \vect e_k\T\mat P_k^{-1}\vect e_k$ with $\vect e_k = \hx_k - \vect x_k$, heading component wrapped.
\end{enumerate}

\section*{Two reminders}
\begin{itemize}[nosep]
  \item The mean always goes through $\vect f_d$, never through $\mat F\hx + \mat G\vect u$.
  \item $\mat F_k$ and $\mat G_k$ are the derivatives of whatever $\vect f_d$ is actually implemented. With the second-order map $\vect p_{k+1} = \vect p_k + T\vect v_k + \tfrac{T^2}{2}\Rot(\theta_k)\vect a_k$ instead of Euler, the position rows of $\mat F_k$ gain $\tfrac{T^2}{2}\,\partial(\Rot\vect a)/\partial\theta$ in the heading column and the position rows of $\mat G_k$ gain $\tfrac{T^2}{2}\Rot(\theta_k)$ in the accelerometer columns.
\end{itemize}

\end{document}
